\section{Related Work}
\label{sec:related}

\textbf{Traditional Computer Vision:} Early approaches to agricultural object detection relied heavily on hand-crafted heuristic features, utilizing algorithms focused on edge detection and color thresholding. While functionally simple, these models struggled with generalization. In complex agricultural environments, bounding features are highly unstable; models frequently failed when confronted with occlusion, target camouflage against green canopies, and variable illumination gradients.

\textbf{Two-Stage Deep Learning Detectors:} The paradigm shifted with the introduction of Convolutional Neural Networks (CNNs). Early heavyweight architectures, such as R-CNN \cite{girshick2014rich}, utilized Selective Search to generate thousands of region proposals per image, applying explicit neural network training to each region. While this established a baseline for mean Average Precision (mAP), the computational overhead was massive. Subsequent iterations, namely Fast R-CNN and Faster R-CNN, optimized this pipeline by sharing convolutional feature maps and introducing the Region Proposal Network (RPN). Despite achieving high localization accuracy, the inherent two-stage architecture (proposal generation followed by classification) remains computationally expensive, failing to meet the strict ultra-low latency requirements for real-time inference on edge devices.

\textbf{Single-Stage Detectors and Vision Transformers:} To achieve real-time detection, single-stage architectures like YOLO (You Only Look Once) \cite{redmon2016you} reframed object detection as a unified regression problem, passing the full image through a single network to simultaneously predict spatial coordinates and class probabilities. While early YOLO versions struggled with small-object detection due to aggressive downsampling, architectures like SSD (Single Shot MultiBox Detector) introduced hierarchical, multi-scale feature maps to improve the receptive field. More recently, attention-based models such as DETR \cite{carion2020end} and the underlying Transformer architecture \cite{vaswani2017attention} have achieved state-of-the-art accuracy; however, the quadratic computational complexity of the transformer attention mechanism makes them prohibitively slow for real-time agricultural deployment.

\textbf{Our Approach:} To deploy an actionable agronomic system, the selected architecture must strike an optimal Pareto frontier between inference velocity and feature extraction accuracy. Transformer-based models are computationally prohibitive for Unmanned Aerial Vehicle (UAV) edge hardware, and two-stage detectors suffer from latency. Therefore, we utilize an advanced iteration of the YOLO architecture, specifically optimized to handle the dense, small-scale nature of the Global Wheat dataset in real-time.